Trabajar con dos \dfnpl{agente} a la vez sobre el mismo repositorio plantea los mismos
problemas que un equipo humano (conflictos al unir ramas, cambios inadvertidos en código
ajeno, decisiones incompatibles) con una diferencia: los agentes no se comunican entre sí
salvo a través de lo que queda escrito en el repositorio. Los mecanismos de esta sección se
diseñaron con esa premisa: toda coordinación es un fichero versionado y, cuando es posible,
una comprobación automática.

\subsection{Aislamiento: ramas, \textit{worktrees} y áreas}

Cada agente trabaja en su propio \dfn{worktree} de Git \cite{git_worktree}:
\texttt{SCDE-claude} sobre la rama \texttt{claude/core-flows} y \texttt{SCDE-codex} sobre
\texttt{codex/admin-ops}, ambas creadas desde la rama de la base común,
\texttt{frontend/base-v2}. Al no compartir directorio de trabajo, ninguno ve los ficheros a
medio editar del otro, y las dependencias, las compilaciones y las pruebas de cada uno son
independientes. Las reglas comunes prohíben a ambos agentes subir cambios al remoto, unir o
reorganizar ramas y reescribir el historial; la integración es una decisión del autor.

El reparto siguió las dependencias del contrato (\cref{TAB:IMP-SPLIT}). Las áreas de la
línea de gestión usan operaciones que el núcleo docente no necesita, de modo que ambas
líneas pudieron avanzar sin esperarse. El único punto de contacto funcional es el enlace de
la página de una asignatura a su pantalla de gestión, que pertenece a la otra línea.

\begin{table}[htbp]
  \centering
  \small
  \begin{tabular}{p{0.44\textwidth}p{0.44\textwidth}}
    \toprule
    \textbf{Núcleo docente (Claude Code)} & \textbf{Gestión (Codex)} \\
    \midrule
    Inicio de sesión, inicio y asignaturas; creación y configuración de exámenes;
    incidencias de escaneo; corrección (\acs{pdf}, anotaciones, notas, \acs{ocr});
    resultados del estudiante; revisiones. Base común y memoria. &
    Usuarios, cursos, configuración y auditoría de la organización; miembros, grupos y
    permisos de cada asignatura; perfil; notificaciones. Publicación de la demostración. \\
    \bottomrule
  \end{tabular}
  \caption{Reparto de áreas entre las dos líneas de trabajo.\label{TAB:IMP-SPLIT}}
\end{table}

\subsection{Propiedad de ficheros y su comprobación}

El fichero \texttt{docs/coordination/OWNERSHIP.json} asigna a cada agente las rutas que
puede modificar (sus directorios de \texttt{src/features/} y sus documentos) y declara como
\emph{compartidas} y congeladas el resto: el sistema de diseño, la infraestructura, el
armazón de la aplicación, la configuración del proyecto, las instrucciones de los agentes y
los documentos de coordinación. La línea del núcleo actúa además como responsable de la base
común, con permiso para modificar las rutas compartidas a condición de anunciar cada cambio.

El \textit{script} \texttt{ownership-check.mjs} compara la rama con la base común y falla si
algún fichero modificado queda fuera de las rutas del agente. En su primera versión solo
examinaba los \textit{commits}; la línea de gestión advirtió que un fichero preparado pero
no consolidado escapaba a la comprobación, y se amplió para revisar también el índice y el
directorio de trabajo. La comprobación forma parte de la condición de cierre de cada
iteración (\cref{SEC:IMP-METHOD}).

\subsection{Peticiones a la base común}

Cuando una línea necesita un cambio en una ruta compartida no lo hace: lo solicita en
\texttt{docs/coordination/REQUESTS.md}, un fichero en el que solo se añaden líneas. La línea
de gestión registró por esta vía tanto huecos del contrato (la falta de un listado de cursos
o de una forma de descubrir la membresía propia) como necesidades de la base: accesibilidad
de los diálogos, los avisos y los campos de formulario, la actualización de la sesión tras
cambiar el correo en el perfil, la ampliación de la auditoría de propiedad y el enlace a la
gestión de asignaturas. Todas las relativas a la base se atendieron desde la línea del
núcleo, que anotó en el mismo fichero cada cambio hecho en la base y su motivo. Al ser de solo
añadidos, los conflictos que produce al unir ramas se resuelven conservando las líneas de
ambas partes.

Cada línea mantiene además su propio registro: el de la línea de gestión
(\texttt{docs/coordination/codex-log.md}) recoge decisiones, errores y órdenes ejecutadas, y
el de la línea del núcleo se integra en el registro técnico del proyecto
(\texttt{docs/tfg/}). Estos registros son la fuente del \cref{AP:ITERATIONS}.

\subsection{Auditoría e integración}

La línea del núcleo audita periódicamente la rama paralela sin modificarla: propiedad de
ficheros, auditoría de estilo y uso exclusivo de operaciones presentes en el inventario de
la \acs{api}. Para ello se definió un subagente de solo lectura
(\texttt{.claude/agents/parallel-auditor.md}). En la primera auditoría, los 25 ficheros
modificados por la línea de gestión estaban dentro de sus rutas, no había errores de estilo
y las 28 operaciones que usaba figuraban en el inventario.

La integración se comprobó en una rama desechable, \texttt{preview/integration}, que une
ambas líneas para ejecutar sobre el conjunto la misma verificación que cada una pasa por
separado. Allí apareció el único fallo de integración del proyecto: una prueba de la línea
de gestión fallaba solo al unir las ramas porque el manejador simulado de
\texttt{GET /subjects/:id/}, registrado antes, capturaba también \texttt{GET
/subjects/export/}. La causa estaba en la base común (el enrutador de la simulación elegía
la primera ruta registrada) y se corrigió allí, eligiendo la ruta más específica como hace
Django, con una prueba de regresión (\cref{SEC:DES-DATA}). En las uniones posteriores no
hubo conflictos en el código; solo en los ficheros de solo añadidos, resueltos por unión.

En la última fase, la publicación de la demostración se asignó a la línea de gestión. Para
no depender del estado del directorio de trabajo de la otra línea, que podía estar a mitad
de una tarea, la composición se construye a partir del último \textit{commit} de cada rama
(\texttt{git archive}), y se comprueba que el contenido publicado coincide con el compilado.

\subsection{Incidencias del trabajo en paralelo}

El aislamiento no fue completo. Al instalar las dependencias en la máquina compartida,
\texttt{npm} modificó el fichero de bloqueo del \dfn{worktree} de la línea de gestión; se
detectó con \texttt{git status} y se revirtió antes de que esa línea empezase a trabajar.
Con los dos agentes, sus servidores de desarrollo y sus pruebas ejecutándose a la vez, varias
pruebas superaron los tiempos de espera por defecto de Testing Library, lo que obligó a
ampliarlos (\cref{AP:ITERATIONS}, entrada 12). Por último, el entorno aislado en que se
ejecutaba Codex impidió en ocasiones lanzar procesos o escribir en el índice de Git, de modo
que algunas de sus iteraciones terminaron con el trabajo verificado pero sin
\textit{commit}; esas iteraciones lo registraron así en su bitácora en lugar de atribuirse
un resultado no consolidado.

\pendiente{Valoración del autor: qué funcionó y qué cambiaría de esta forma de coordinar
dos agentes, y cuánto tiempo de supervisión le exigió.}
